% from ega4-21c-fr.tex
\oldpage[IV-4]{317}
Si en outre $A$ est noethérien, la condition (i) équivaut à
\begin{enumerate}
  \item[\textup{(i bis)}] $\operatorname{prof}(A)\geq2$.
\end{enumerate}
\end{proposition}

\begin{proof}
En effet, le seul ouvert de $X$ contenant $a$ est $X$ lui-même, et par suite tout
$\mathcal O_X$-Module inversible est isomorphe à $\mathcal O_X$, autrement dit
$\operatorname{Pic}(X)=0$; la première assertion résulte donc de la définition
(21.13.1) et de (21.13.3). La seconde assertion n'est autre qu'un cas particulier de
(21.13.4).
\end{proof}

\begin{env}[21.13.9]
\label{IV.21.13.9}
\emph{Exemples}. ---

(i) Un anneau local noethérien parafactoriel est nécessairement de dimension
$\geq2$ en vertu de (21.13.8); en d'autres termes un anneau local noethérien de
dimension $\leq1$ n'est pas parafactoriel.

(ii) Un anneau local noethérien \emph{factoriel} $A$ de dimension $\geq2$ est
parafactoriel, comme il résulte de (21.13.8) et de (21.6.14).

(iii) Si $A$ est un anneau local noethérien de dimension $\geq3$ et parafactoriel, il
n'est pas nécessairement normal ni même réduit. Considérons un anneau local régulier
$B$ de dimension $\geq3$, et soit $A=D_B(B)$ $(\mathbf 0, 18.2.3)$, isomorphe à
$B[T]/(T^2)$; on voit aussitôt que l'on a
$\operatorname{prof}(A)=\operatorname{prof}(B)=\dim(B)\geq3$. Pour montrer que $A$
est parafactoriel, il suffit donc, avec les notations de (21.13.8), de prouver que
$\operatorname{Pic}(U)=0$. Soit $\mathfrak J$ le noyau de l'augmentation $A\to B$,
qui est tel que $\mathfrak J^2=0$ et qui, en tant que $B$-module, est isomorphe à $B$.
Comme $B=A/\mathfrak J$, $X=\operatorname{Spec}(A)$ et
$X_0=\operatorname{Spec}(B)$ ont même espace sous-jacent; si
$\mathscr J=\widetilde{\mathfrak J}$, $X_0$ est le sous-schéma défini par
$\mathscr J$; nous noterons $U_0$ le sous-préschéma induit par $X_0$ sur l'ouvert
$U$. Pour tout $z\in\mathfrak J$, posons $\varphi(z)=1+z$; puisque
$(1+z)(1-z)=1$, $1+z$ est inversible dans $A$ et $\varphi(\mathfrak J)$ est le noyau
de l'homomorphisme surjectif canonique de groupes multiplicatifs $A^*\to B^*$; en
d'autres termes, on a une suite exacte de groupes commutatifs
\[
  0\longrightarrow\mathfrak J\xrightarrow{\varphi}A^*\longrightarrow B^*
  \longrightarrow1
\]
(les trois derniers groupes étant multiplicatifs, les deux premiers additifs). Pour la
même raison, pour tout $t\in A$, on a la suite exacte
\[
  0\longrightarrow\mathfrak J_t\xrightarrow{\varphi_t}(A_t)^*
  \longrightarrow(B_{t_0})^*\longrightarrow1
\]
en désignant par $t_0$ l'image de $t$ dans $B$, puisque
$B_{t_0}=A_t/\mathfrak J_t$; autrement dit, on a une suite exacte de faisceaux de
groupes commutatifs sur l'espace topologique $X$
\[
  0\longrightarrow\mathscr J\xrightarrow{u}\mathcal O_X^*
  \longrightarrow\mathcal O_{X_0}^*\longrightarrow1
\]
d'où par restriction à l'ouvert $U$, une suite exacte
\[
  0\longrightarrow\mathscr J|U\longrightarrow\mathcal O_U^*
  \longrightarrow\mathcal O_{U_0}^*\longrightarrow1.
\]
Par la suite exacte de cohomologie, on en déduit une suite exacte
\[
\label{IV.21.13.9.1}
  \mathrm H^1(U,\mathscr J)\longrightarrow\mathrm H^1(U,\mathcal O_U^*)
  \xrightarrow{u_1}\mathrm H^1(U,\mathcal O_{U_0}^*)
  \tag{21.13.9.1}
\]
Mais puisque $\mathfrak J$ est un $B$-module isomorphe à $B$, on a
$\mathrm H^1(U,\mathscr J)=\mathrm H^1(U_0,\mathcal O_{U_0})$, et il résulte du
chap.~III, 3\textsuperscript{e} partie (voir aussi [41, III, Exemple III-1]) que l'on a
